\section{Related Work}
\label{sec:related}

Our study sits at the intersection of three lines of research: (i) empirical
work characterizing what makes bug reports useful, (ii) techniques that assess
and improve the \emph{steps to reproduce} (S2R) that determine whether a report
is reproducible, and (iii) the emerging use of large language models (LLMs) for
bug report analysis and as automated evaluators. We review each theme and then
position our contribution.

\subsection{Bug Report Quality and Information Completeness}
\label{sec:related:quality}

A long line of empirical research establishes that bug reports vary widely in
quality and that this variation directly affects triage and resolution effort.
Survey evidence from developers of large open-source projects shows a systematic
mismatch between the information developers most need---steps to reproduce,
stack traces, and test cases---and the information reporters typically
provide~\cite{zimmermann2010good}. Complementary quantitative studies model
report quality from surface and textual features and confirm that such features
predict how quickly a report is acted upon, framing quality assessment as a
largely automatable task~\cite{hooimeijer2007modeling}. Broader content analyses
across multiple projects further quantify how often key fields are missing and
trace the downstream consequences of incomplete reports for developers who must
recover the missing context~\cite{davies2014whats}. Collectively, this theme
motivates \emph{reproducibility} as the quality dimension of greatest practical
value, yet it treats quality mostly through proxy features rather than by
judging whether a report is actually reproducible.

\subsection{Reproducibility and the Quality of Steps to Reproduce}
\label{sec:related:s2r}

A second theme targets the S2R directly, since incomplete or ambiguous steps are
the dominant obstacle to reproducing a failure. Early tool support augmented the
reporting process itself, guiding Android reporters to produce more complete,
reproducible steps at submission time~\cite{moran2015autocompleting}. Later work
shifted toward automatically \emph{assessing} S2R quality: Euler identifies the
reported steps and flags missing or low-quality ones to give reporters
actionable feedback~\cite{chaparro2019assessing}, and more recent approaches
combine language understanding with app UI models so that natural-language steps
are grounded in concrete interface interactions~\cite{mahmud2025combining}.
Closest to our setting, ImproBR uses an LLM pipeline to detect and repair
missing or ambiguous S2R, observed behavior, and expected behavior on the Mojira
issue tracker, substantially raising the share of executable steps and fully
reproducible reports~\cite{akyol2026improbr}. A common thread, however, is that
these S2R techniques are built for GUI/mobile applications and rely on program
or dynamic analysis of a runnable app to establish the language-to-program
mapping---an assumption that does not transfer to a game issue tracker such as
Minecraft/Mojira, and none of them frames reproducibility assessment as a
measured agreement task against human judgment.

\subsection{Large Language Models for Bug Report Analysis and Evaluation}
\label{sec:related:llm}

The third theme concerns LLMs applied to software engineering and, increasingly,
as evaluators. Systematic reviews document the rapid adoption of LLMs across SE
tasks and note that data quality and rigorous evaluation---not merely model or
dataset scale---determine reliability~\cite{hou2024llm4se}. In parallel, the
``LLM-as-judge'' paradigm studies when a model's assessments can substitute for
human evaluation, together with the meta-evaluation methods and limitations that
govern their trustworthiness~\cite{li2024judges}. This literature makes clear
that an LLM used as an assessor must itself be validated against human
annotators using an explicit reliability metric, rather than assumed correct.

\subsection{Positioning of Our Work}
\label{sec:related:positioning}

Prior work either (i) characterizes bug report quality through proxy features
without judging actual reproducibility, (ii) assesses S2R quality for GUI/mobile
apps using program or dynamic analysis that is infeasible for the
Minecraft/Mojira domain, or (iii) applies LLMs to \emph{generate}, \emph{improve},
or \emph{classify} reports---including ImproBR on the same Mojira
data~\cite{akyol2026improbr}---without validating the LLM as a reproducibility
\emph{assessor}. Our study addresses this gap: we use an LLM to automatically
assess bug report reproducibility on Mojira and evaluate it directly against
human expert annotation, treating the LLM-as-judge concern
seriously~\cite{li2024judges} by reporting Cohen's Kappa against a pre-registered
threshold of $\kappa \geq 0.70$. Uniquely, we run this evaluation on \emph{both}
the original reports and the LLM-improved versions that ImproBR produces from
them~\cite{akyol2026improbr}, treating the two as co-equal conditions, which lets
us measure whether an LLM evaluator's agreement with humans holds up on text that
another LLM has rewritten. To our knowledge, this is the first study to benchmark
automated reproducibility assessment as a human-agreement task on this domain,
and the first to quantify how that agreement changes on AI-improved reports.
